\chapterfont{\color{Grey}}  % sets colour of chapter

\sectionfont{\color{Grey}}  % sets colour of sections

\subsectionfont{\color{Grey}}  % sets colour of subsections


\renewcommand\pcolor{Grey}
\renewcommand{\headrule}{\hbox to\headwidth{%

		\color{Grey}\leaders\hrule height \headrulewidth\hfill}} % color of title

\fancyfoot[LE,RO]{\thepage}


{ \Large \leftwatermark{

		\put(-76.5,-75){\includegraphics[scale=0.8]{img/thumbindex/thumbindex_Grey.eps}} 

		\put(-67,-66.5){ {\color{white} 1 }}

		\put(-67,-91.5){ 2 }

		\put(-67,-116.5){ 3 }

		\put(-67,-141.5){ 4 }

		\put(-67,-166.5){ 5 }

		\put(-67,-191.5){ 6 }

	} \rightwatermark{

		\put(346.5,-75){\includegraphics[scale=0.8]{img/thumbindex/thumbindex_Grey.eps}}  \put(350.5,-66.5){ {\color{white} 1 }}

		\put(350.5,-91.5){ 2 }

		\put(350.5,-116.5){ 3 }

		\put(350.5,-141.5){ 4 }

		\put(350.5,-166.5){ 5 }

		\put(350.5,-191.5){ 6 }

}}



\chapter{Introduction}

\label{chap:introduction}



\newpage


\noindent


When the first draft of the human genome was released in June of 2000, the then president of the United States, Bill Clinton, told the press that the genome sequence would "revolutionize the diagnosis, prevention, and treatment of most, if not all, human diseases"\cite{collinsHasRevolutionArrived2010}. Without a doubt, the first human genome sequence has revolutionized genomic research. Until that time, genetic studies had mostly focused on rare, monogenic diseases using methods like linkage analysis and Sanger sequencing. At this point there were 2000 genes known that were causal for these rare monogenic diseases. In the 18 years after the first draft of the human genome was completed, there was a four-fold increase in the number of causal genes identified\cite{claussnitzerBriefHistoryHuman2020}. In addition, further developments of computational and DNA sequencing technologies invented or optimized for the human genome project\cite{hoodHumanGenomeProject2013} have helped to bring down the costs of DNA sequencing dramatically, from over \$2 billion for the first human genome to less than \$1000 now. This reduction in cost enabled large projects like the 1000 Genomes project\cite{the1000genomesprojectconsortiumGlobalReferenceHuman2015} or the Genome of the Netherlands \cite{ fokkemaDutchGenomeDiagnostic2019} and ENCODE\cite{consortiumENCODEENCyclopediaDNA2004} to produce maps of sequence variation and to identify functional elements in the human genome. 


The impact of the project was, and still is, enormous and has enabled many advances within genetics and genomics research\cite{chicheBenchtobedsideReviewFulfilling2002} and in identification of new disease genes and has helped to improve genetic diagnosis\cite{claussnitzerBriefHistoryHuman2020}. However, this has only led to treatment in a limited number of cases. Even though diagnosis of genetic disorders has improved, genetic testing is still unable to find the causal genes and variants in the majority of cases\cite{diemenRapidTargetedGenomics2017}. Genetic tests are now successfully used for preventative care of monogenic disorders, such as preventative mastectomies or extra surveillance for BRCA1/2 carriers \cite{heemskerk-gerritsenSurvivalBilateralRiskreducing2019}, but for complex disorders, where many variants have a small effect as opposed to one or a few variants having a large effect, this is still difficult\cite{claussnitzerBriefHistoryHuman2020}. 


Despite the fact that the past 20 years have been full of incredible advances in genomics and genetics research, we are not yet close to the prevention and treatment of most human diseases. Still, there are many ways in which the field of genetics can assist in improving clinical outcome for rare and complex disorders. Currently, there are public-private partnerships under way that will help better our understanding of genotype-phenotype relationships\cite{HomepageInternationalCommon}, provide large data resources for researchers\cite{fingennFinnGenDocumentationR32020,sudlowUKBiobankOpen2015}, and assist in drug target identification\cite{carvalho-silvaOpenTargetsPlatform2019}. Improvements in polygenic risk scores will allow prediction of risk for polygenic traits\cite{natarajanpradeepPolygenicRiskScore2017, kheraGenomewidePolygenicScores2018}, which can help in preventative care for complex disorders. In drug development, the fraction of drug compounds with direct genetic support increases substantially between the preclinical and approved stage (2.0\% versus 8.2\%, respectively\cite{nelsonSupportHumanGenetic2015}), which demonstrates the added benefit of integrating genetic research into the drug development pipeline. %Finally, integrating layers of omics and environmental data will address the issue of “missing heritability”\cite{claussnitzerBriefHistoryHuman2020}. 


In this thesis, I address two important factors that are needed if we want to use functional genomics to assist in improving healthcare: (1) insight into the effects of rare variants on gene expression and the genetic architecture of molecular traits and (2) the cell-type- and tissue-context of the genetic effects on molecular traits.


\section{General introduction}

\subsection{Gene regulation through genetic variation}

The central dogma of molecular biology states that genes, the genetic information stored in DNA, are transcribed to RNA, which is then translated to a protein\cite{crickCentralDogmaMolecular1970}. Almost every process in cells is then regulated by these proteins. They act as transcription factors (regulators of gene activity), antibodies (binding to specific foreign particles), enzymes (enablers of many chemical reactions), messengers, and structural components, as well as having roles in transport or storage\cite{uzmanMolecularBiologyCell2003}. However, exceptions to this dogma exist: there are many genes for which the RNA transcript does not get translated to a protein, yet these non-coding genes are often involved in regulating the activity of genes\cite{shabalinaMammalianTranscriptomeFunction2004}. 


The DNA that contains the genetic information consists of about 6.4 billion nucleotide bases divided over 23 chromosome pairs. When comparing a typical human genome sequence to the sequence of the reference human genome, about 4.1 to 5 million sites will be different (< 0.1\%). Most often, these differences are either a difference of a single nucleotide or the insertion or deletion of a single base\cite{the1000genomesprojectconsortiumGlobalReferenceHuman2015}. Large deletions, duplications, inversions, and other chromosome changes occur, but the focus of this thesis is on single nucleotide changes and short insertions and deletions. Many of these differences will have no noticeable effect on any observable characteristics (the phenotype). Others will have a subtle effect, e.g. having a very small effect on adult height, of which some can also increase risk for common diseases. A smaller fraction of these variants can more profound effects on certain phenotypes, for example eye colour. And sometime, other rare variants can cause rare but severe genetic disorder such as cystic fibrosis or Huntington’s disease. 


\subsection{The effects of single nucleotide polymorphisms}

Single Nucleotide Polymorphisms (SNPs), the difference of a single nucleotide base, can affect a gene in multiple ways. A SNP within a gene (coding region) can cause the protein product to be prematurely stopped\cite{deboeverMedicalRelevanceProteintruncating2018} or to fold differently\cite{vihinenTypesEffectsProtein2015}, thus reducing or completely removing its functionality. A SNP can also be located outside of the gene (in a non-coding region), for example in the start site of a gene or in an enhancer region, and thereby affecting gene activity. A SNP at a transcription start site can cause a transcription factor to bind with less affinity, thereby lowering gene activity, which can also lead to disorders. One example of this is the Pierre Robin sequence, which is caused by changes in non-coding sequences around the SOX9 gene\cite{benkoHighlyConservedNoncoding2009}. However, variation in a non-coding region most often has a smaller effect than mutations in a coding region, and it usually takes multiple non-coding mutations to cause a severe disorder.


\subsection{Genetic disorders}

Genetic disorders are often divided in two broad classes: monogenic disorders\cite{heidiRareGeneticDisorders2008} and complex disorders\cite{craigComplexDiseasesResearch2008}. Monogenic disorders are caused by large-effect variants where a SNP in a single gene is enough to cause a disease. Some well-known examples of monogenic disorders are cystic fibrosis\cite{keremIdentificationCysticFibrosis1989}, sickle-cell anaemia\cite{neelInheritanceSickleCell1949}, and Huntington's disease\cite{gusellaPolymorphicDNAMarker1983}. Complex disorders, on the other hand, are characterised by multiple variants, usually in regulatory regions, that affect the expression levels of multiple genes. Each of these complex disease SNPs confer a limited amount of risk for getting the disease. Examples of complex disorders include diabetes mellitus type 2\cite{americandiabetesassociationDiagnosisClassificationDiabetes2007}, Alzheimer's disease\cite{hardyAlzheimerDiseaseAmyloid1992}, and autism\cite{PolygenicTransmissionDisequilibrium}. For monogenic disorders, it is relatively easy to identify the causal variant because they are most often located within genes and either change important amino acids or cause frame shifts. Identifying associated variants for complex disorders is, however, much more difficult because each variant only has a small contribution, and environmental factors can play an important role as well. Once associated variants have been found, the next challenge: to identify the causal variant remains difficult. For common variants often numerous variants exists that are in strong linkage disequilibrium, making it very challenging to pinpoint the true causal variant.


One step towards identifying associated variants is to create resources of genotype‒phenotype relationships that can be easily used by many researchers.\cite{claussnitzerBriefHistoryHuman2020}. Thanks to large collaborative efforts, there are now biobanks with genotype data for more than 100,000 individuals (for example, the UK biobank\cite{sudlowUKBiobankOpen2015}, LifeLines\cite{tigchelaarCohortProfileLifeLines2015}, FinnGenn\cite{fingennFinnGenDocumentationR32020} and the Estonian Genome Project\cite{metspaluEstonianGenomeProject2004}). With these large datasets, it is possible to identify which DNA sequence variants are associated with increased or decreased risk for many different (disease) phenotypes at once, by comparing the genotypes of a large number of healthy people with those of a large group of people with a specific disease or other phenotype. These studies are called Genome Wide Association Studies (GWAS), and many of them are now collected in the GWAS Catalog\cite{bunielloNHGRIEBIGWASCatalog2019}, with the 14 July 2020 GWAS Catalog release containing 116,480 variants for 4,467 phenotypes. The question then remains: What are the molecular mechanisms in which these disease-associated genes contribute to disease risk? One way to discern this is by mapping quantitative trait loci\cite{membersofthecomplextraitconsortiumNatureIdentificationQuantitative2003} (QTLs). Quantitative traits can be any measurable phenotype that varies in the human population due to genetics, e.g. height or molecular traits such as gene expression, protein, or methylation levels. A QTL is a locus, a region on the DNA, where variation affects the quantitative trait. Disease-associated SNPs from GWAS studies can be linked to changes in, for example, expression levels, by mapping expression QTLs (eQTLs) for all GWAS SNPs.


The effects that genetic variants have on gene expression can be different depending on the tissue and cell type (\textbf{Figure 1}), therefore these genotype-phenotype maps have to be made for all tissues. 

%<add this somewhere maybe Natural Selection on Genes that Underlie Human Disease Susceptibility https://www-sciencedirect-com.proxy-ub.rug.nl/science/article/pii/S0960982208006015\cite{blekhmanNaturalSelectionGenes2008}> 


%\subsection{Human genetic evidence for drug development}

%One area where genetics can be of definite help is in the pre-selection of candidate drug compounds. It is currently very costly to perform clinical drug trials, and, even though there have been many technological advances in recent decades, drug development costs are high while there has been little increase in the output of new drugs \cite{cookLessonsLearnedFate2014a}. By including genetic information in drug compound selection, the success rate in clinical development can be doubled\cite{nelsonSupportHumanGenetic2015b}. 


\section{This thesis}

The aim of this thesis is to get a better understanding of the context in which genetic risk factors influence gene expression levels. In Part 1, we describe the effects of rare variants on gene expression in blood and the genetic architecture of molecular traits. In Part 2, we identify genetic regulation in the context of different cell types in blood and brain.


\subsection{Part 1 - Architecture of molecular traits and rare variant signals}

In Chapter 2, we compared the heritability explained by all the studies included in the GWAS Catalog\cite{bunielloNHGRIEBIGWASCatalog2019} to the heritability explained by expression quantitative trait loci (eQTL) and methylation quantitative trait loci (meQTL) studies. By doing so, we find that the genetic architecture of molecular traits is less polygenic than that of (disease) phenotypes. In chapter 3, we identify the disease categories for which allele-specific expression can be informative for prioritizing putative causal variants.


\subsection{Part 2 - Cell type- and tissue-context of genetic regulation of gene expression}

In chapter 4, we develop two methods. The first predicts the cell count proportions from whole blood samples using gene expression data. The second deconvolutes the cell type‒interacting eQTL effects found in whole blood samples. In chapter 5, we compare the genetic regulation of gene expression in brain and blood and identify which genes are affected by neurodegenerative and psychiatric risk factors. In chapter 6, I discuss how these findings can contribute to improving health care outcomes.




\bibliographystyle{naturemag}

\bibliography{chapters/chapter1-introduction/chapter1-introduction}



% \cite{gibbsHumanGenomeProject2020} It should be noted that the hype was from front-facing... <- read this piece 

